% 示例阅读顺序：查看实际效果，在源码中复制对应 frame，再按需替换内容。
\section{文字与列表}
\subsection{页面与文字}

% 含 verbatim 的 frame 必须使用 fragile；示例中的结束行带注释，避免被外层提前识别。
\begin{frame}[fragile]{页面与章节的基本写法}
  \begin{verbatim}
\section{章节名}
\subsection{小节名}
\begin{frame}{页面标题}
  页面正文。
\end{frame} % 结束示例页面
  \end{verbatim}
  \slideNote{章节和小节会自动生成过渡页。星号形式不进入目录，
  也不触发过渡页；封面标题与日期在 config/metadata.tex 中设置。}
\end{frame}

\begin{frame}{文字、强调与页面副标题}
  \framesubtitle{副标题使用 framesubtitle 命令设置}
  正文支持中文与 English 混排，段落之间使用空行分隔。

  \textbf{粗体}用于强调；\textit{Italic text}展示西文斜体；
  \texttt{inline code}用于代码片段。

  \keypoint{主题色粗体重点}，以及 \alert{需要强调的条件}。
  \slideNote{正文保留黑体；页标题和表格使用宋体。注释文字自动使用小一号字号，
  与正文保持清晰的层级。}
\end{frame}

\subsection{列表与代码}
% 嵌套列表：在 item 内再放入 itemize / enumerate 环境。
\begin{frame}{无序、有序与嵌套列表}
  \begin{columns}[T,onlytextwidth]
    \column{0.48\textwidth}
      \textbf{无序列表：itemize}
      \begin{itemize}
        \item 第一项。
          \begin{itemize}
            \item 嵌套子项。
            \item 另一条子项。
          \end{itemize}
        \item 第二项。
      \end{itemize}
    \column{0.48\textwidth}
      \textbf{有序列表：enumerate}
      \begin{enumerate}
        \item 第一步。
          \begin{enumerate}
            \item 第一步的子步骤。
          \end{enumerate}
        \item 第二步。
      \end{enumerate}
  \end{columns}
  \slideNote{使用 item 添加条目；在列表内部再放一个列表环境即可嵌套。}
\end{frame}

\begin{frame}{宽松要点与术语列表}
  \begin{columns}[T,onlytextwidth]
    \column{0.48\textwidth}
      \textbf{keypoints：宽松要点}
      \begin{keypoints}
        \item 第一个核心观点。
        \item 第二个核心观点。
        \item 第三个核心观点。
      \end{keypoints}
    \column{0.48\textwidth}
      \textbf{description：术语说明}
      \begin{description}[模型]
        \item[数据] 输入与观测。
        \item[模型] 结构与参数。
        \item[指标] 评价的依据。
      \end{description}
  \end{columns}
\end{frame}

\begin{frame}{网址、路径与特殊字符}
  \begin{itemize}
    \item 网址：\url{https://ctan.org/pkg/beamer}
    \item 命名链接：\href{https://ctan.org/pkg/beamer}{Beamer 文档入口}
    \item 路径：\path{figures/example-figure.png}
    \item 行内代码：\texttt{tectonic main.tex}
  \end{itemize}

  \begin{block}{特殊字符需要转义}
    百分号写 \texttt{\textbackslash\%}，与号写 \texttt{\textbackslash\&}，
    下划线写 \texttt{\textbackslash\_}，井号写 \texttt{\textbackslash\#}。
    需要保留原始排版的整段代码，可把 frame 声明为
    \texttt{\textbackslash begin\{frame\}[fragile]} 后使用 verbatim 环境。
  \end{block}
\end{frame}

\begin{frame}[fragile]{整段代码与原样文本}
  页面声明：\verb|\begin{frame}[fragile]{代码示例}|。
  \begin{verbatim}
score = mean(values)
# keep spaces and symbols: _ % &
  \end{verbatim}
  \slideNote{verbatim 保留空格和特殊字符；含有原样代码的页面必须使用 fragile。
  不需要额外的语法高亮宏包或外部程序。}
\end{frame}
